% =====================================================================
% Application draft H2 — Vol I main.tex Part I opener Climax block.
%
% Source: wave_supervisory_climax_main_tex_drafts.md §2 (V22).
% Targets:
%   H2a — INSERT after L917 (closing of "categorical logarithm" sentence)
%         and before existing L919 ("A logarithm has four properties...").
%   H2b — REPLACE existing L919 sentence
%         "A logarithm has four properties. Each is a theorem."
%         with the patched cadence below.
%
% Posture: rhetorical promotion within the proved core.  The FM-residue
% paragraph at L902–916 already builds momentum toward Arnold's three-term
% identity; the Climax block below names the universal pullback that
% all four properties (A)–(D)+(H) below are unfoldings of.
%
% Conventions verified:
%   - \nabla_{Arnold} = d - sum_{i<j} t_{ij} d log(z_i - z_j),
%     with t_{ij} the Drinfeld–Kohno–Bar-Natan infinitesimal braid
%     generators (NO bare scalar prefactor: t_{ij} carries its own level
%     when specialised, AP126/AP141).
%   - Cross-reference Part~\ref{part:characteristic-datum} (live label,
%     main.tex L1096).
%   - bare \kappa is Vol I house style.
%   - Connective patch (H2b) inserts six words: "proved by such an
%     identity"; "structural"; "; each is one face of the Climax",
%     preserving the existing four-property enumeration verbatim.
%
% No new macros required: \cA, \Conf, \Einf, \barB, \fg, \CC, \hv,
% \ClaimStatus*, \ChirHoch all live in main.tex preamble.
%
% Status of forward references at insertion time:
%   - \ref{part:characteristic-datum}  EXISTS.
%   - The Climax equations are STATED, not cited as theorems; the
%     numbered theorem environment thm:climax is left for a follow-on
%     commit (Obstruction §7.1 of the supervisory report).
% =====================================================================


% ----- H2a: insert after L917 (after "categorical logarithm of $\cA$.") -----

\medskip\noindent
\textbf{The Climax of Volume~I.}
The bar differential just constructed by FM residue extraction on
$\overline{C}_{n}(X)$ is the pullback of Arnold's universal KZ
connection on $\Conf_{n}^{\mathrm{ord}}(X)$:
\[
  d_{\mathrm{bar}}
  \;=\;
  \mathrm{KZ}^{*}(\nabla_{\mathrm{Arnold}}),
  \qquad
  \nabla_{\mathrm{Arnold}}
  \;=\;
  d \;-\; \sum_{i<j} t_{ij}\;d\log(z_{i}-z_{j}),
\]
with $\{t_{ij}\}$ the universal infinitesimal braid generators
(Arnold 1969). Flatness $\nabla_{\mathrm{Arnold}}^{2} = 0$ is
equivalent, after pullback, to $d_{\mathrm{bar}}^{2} = 0$;
Arnold's three-term identity on $\Conf_{3}^{\mathrm{ord}}(\CC)$
is the chain-level expression of the classical Yang--Baxter equation
on the collision residues
$r^{(ij)}(z_{i}-z_{j}) = \mathrm{KZ}^{*}(\eta_{ij})$.

Three classical theorems --- Drinfeld--Kohno
$\rho_{n}^{\mathrm{KZ}} \cong \rho_{n}^{U_{q}(\fg)}$,
Verlinde $N_{ij}^{k} = \sum_{a} S_{ia}S_{ja}\overline{S_{ka}}/S_{0a}$,
and the Borcherds product $\Phi_{10}$ --- are pullbacks of
$\nabla_{\mathrm{Arnold}}$ along three structure functors,
applied to three specific chiral inputs (affine Kac--Moody on
evaluation modules; rational chiral algebra on conformal blocks;
lattice vertex algebra on the genus-$1$ trace). The bar
differential is a fourth pullback --- the chain-level
incarnation. The four properties (A)--(D) below, and the
modular-characteristic identity
$\kappa(\cA) = -c_{\mathrm{ghost}}(\mathrm{BRST}(\cA))$ established
in Part~\ref{part:characteristic-datum}, are the structural
unfolding of this single observation.


% ----- H2b: replace the existing L919 sentence
%   OLD:  A logarithm has four properties. Each is a theorem.
%   NEW (one line, preserving the connective cadence):

\smallskip\noindent
A logarithm proved by such an identity has four structural
properties. Each is a theorem; each is one face of the Climax.
